% 对应题号 M1-55；图与本题条件、推导、结论逐项对应。
% 图契约：印刷宽16cm；最小字号8.4pt；几何/场或曲线在左，计算关系在右。
\documentclass[tikz,border=7pt]{standalone}
\usepackage{fontspec,xeCJK,amsmath,amssymb,graphicx}
\IfFontExistsTF{Times New Roman}{\setmainfont{Times New Roman}}{\setmainfont{TeX Gyre Termes}}
\IfFontExistsTF{Heiti SC}{\setCJKmainfont[AutoFakeBold=1.5]{Heiti SC}}{\setCJKmainfont[AutoFakeBold=1.5]{FandolHei-Regular}}
\usetikzlibrary{arrows.meta,bending,calc,patterns,positioning}
\definecolor{ink}{HTML}{243745}
\definecolor{blue}{HTML}{275DAB}
\definecolor{green}{HTML}{177C63}
\definecolor{red}{HTML}{BD4935}
\definecolor{orange}{HTML}{AA6C16}
\definecolor{purple}{HTML}{715393}
\tikzset{every node/.style={font=\fontsize{9.3}{12.5}\selectfont,text=ink},
 axis/.style={-{Stealth[length=4pt]},line width=.65pt,draw=ink!65},
 vec/.style={-{Stealth[length=5pt]},line width=1pt,draw=blue},
 short/.style={-{Stealth[length=3pt]},line width=.8pt},
 guide/.style={dashed,draw=ink!40,line width=.5pt},
 card/.style={rounded corners=2pt,draw=blue!65,fill=blue!4,align=center,text width=6.6cm,inner xsep=9pt,inner ysep=7pt},
 title/.style={font=\fontsize{10.5}{14}\selectfont\bfseries,text width=15cm,align=center},
 label/.style={font=\fontsize{8.4}{11}\selectfont,align=center}}
\begin{document}
\begin{tikzpicture}[x=1cm,y=1cm]
\path[use as bounding box] (-8,-3.4) rectangle(8,3.4);
\node[title] at (0,3) {M1-55\quad 边界为零时，调和解的梯度只能处处为零};
\begin{scope}[shift={(-4,0)}]

\filldraw[fill=blue!3,draw=ink!65] (0,0) ellipse(2.3 and 1.7);
\node[blue] at (0,-1.1) {$V:\ \nabla^2u=0$};
\node[label] at (0,2.15) {$\partial V:\ u=0$};
\draw[vec,orange] (-1,-.6)--(.9,.9) node[midway,right] {$\gamma$};
\fill[ink] (-1,-.6) circle(2pt);\fill[ink] (.9,.9) circle(2pt);
\node[label] at (0,-2.6) {区域连通，任意两点可沿区域内路径比较函数值};

\end{scope}\node[card,draw=blue!65,fill=blue!4] at (3.95,1.65) {$\displaystyle\int_V|\nabla u|^2dV=0$};
\node[card,draw=green!65,fill=green!4] at (3.95,0) {被积函数非负且连续，故 $\nabla u=0$。};
\node[card,draw=orange!65,fill=orange!4] at (3.95,-1.65) {$u$ 在连通区域内为常数；\\再由 $u|_{\partial V}=0$ 得 $u\equiv0$。};
\end{tikzpicture}\end{document}
